\documentclass{article}
\usepackage{amsmath, amssymb, amsthm}
\usepackage{booktabs}
\usepackage{multirow}
\usepackage{geometry}
\usepackage{hyperref}
\usepackage{cleveref}

\geometry{margin=1in}
\hypersetup{colorlinks=true, urlcolor=blue}

\newcommand{\N}{\mathbb{N}}
\newcommand{\R}{\mathbb{R}}
\newcommand{\Z}{\mathbb{Z}}
\newcommand{\E}{\mathbb{E}}
\newcommand{\Var}{\mathrm{Var}}
\newcommand{\Prob}{\mathbb{P}}
\newcommand{\Ind}{\mathbf{1}}
\newcommand{\ceil}[1]{\lceil #1 \rceil}
\newcommand{\floor}[1]{\lfloor #1 \rfloor}

\title{Multi-Mode Mixed-Criticality Scheduling}
\author{}
\date{}

\begin{document}

\maketitle

\tableofcontents

\section{Task Model}

See the task model document for details on the formal task model, including:
\begin{itemize}
    \item Basic task properties and criticality assignment
    \item Multi-level degradation model
    \item Job execution model
    \item Schedulability analysis (RTA and AMC-rtb tests)
\end{itemize}

\documentclass{article}
\usepackage{amsmath, amssymb, amsthm}
\usepackage{booktabs}
\usepackage{multirow}
\usepackage{geometry}
\usepackage{hyperref}
\usepackage{cleveref}

\geometry{margin=1in}
\hypersetup{colorlinks=true, urlcolor=blue}

% Shared macros for the multi-mode mixed-criticality scheduling papers.
% % Shared macros for the multi-mode mixed-criticality scheduling papers.
% \input{macros} from any document in this directory rather than redefining
% these locally -- a macro defined in one paper's preamble and silently
% missing from a sibling's is exactly the class of error that put \round and
% \E undefined in evaluation_methodology.tex.

\newcommand{\N}{\mathbb{N}}
\newcommand{\R}{\mathbb{R}}
\newcommand{\Z}{\mathbb{Z}}
\newcommand{\E}{\mathbb{E}}
\newcommand{\Var}{\mathrm{Var}}
\newcommand{\Prob}{\mathbb{P}}
\newcommand{\Ind}{\mathbf{1}}
\newcommand{\ceil}[1]{\lceil #1 \rceil}
\newcommand{\floor}[1]{\lfloor #1 \rfloor}
\newcommand{\round}[1]{\lfloor #1 \rceil}
 from any document in this directory rather than redefining
% these locally -- a macro defined in one paper's preamble and silently
% missing from a sibling's is exactly the class of error that put \round and
% \E undefined in evaluation_methodology.tex.

\newcommand{\N}{\mathbb{N}}
\newcommand{\R}{\mathbb{R}}
\newcommand{\Z}{\mathbb{Z}}
\newcommand{\E}{\mathbb{E}}
\newcommand{\Var}{\mathrm{Var}}
\newcommand{\Prob}{\mathbb{P}}
\newcommand{\Ind}{\mathbf{1}}
\newcommand{\ceil}[1]{\lceil #1 \rceil}
\newcommand{\floor}[1]{\lfloor #1 \rfloor}
\newcommand{\round}[1]{\lfloor #1 \rceil}


\title{Task Model for Multi-Mode Mixed-Criticality Scheduling}
\author{}
\date{}

\begin{document}

\maketitle

\section{Introduction}

This document presents the formal task model used in our research on multi-mode
mixed-criticality (MC) scheduling. The model extends traditional two-mode
Adaptive Mixed-Criticality (AMC) scheduling to support multiple criticality
levels, enabling graded degradation rather than a single normal/degraded
transition.

\textbf{Key analysis principle.} Every degradation level has its own execution
budget, not just normal and fully-degraded. What makes the scheme analysable is
not that only two budgets exist -- they do not -- but that each level's
worst-case execution time (WCET) budget is converted, via the \emph{same}
response-time analysis used for classic two-level AMC, into a worst-case
response time (WCRT) that is monitored at run time. Reaching that WCRT is what
triggers the transition into that level. Safety is preserved because each
level's operating point -- its budget together with the set of LO-criticality
tasks it keeps running -- is checked at design time against exactly this
analysis, generalised across levels rather than restricted to two
(\Cref{sec:admissibility}).

The model builds upon the foundational work presented in~\cite{9804600}.

\section{Task Model}

\subsection{Basic Task Properties}

A task set $\tau = \{\tau_1, \tau_2, \dots, \tau_n\}$ consists of $n$
independent, sporadically arriving tasks. Each task $\tau_i$ is characterized
by:

\begin{table}[h]
\centering
\caption{Task Properties}
\begin{tabular}{lll}
\toprule
Symbol & Name & Description \\
\midrule
$T_i \in \N^+$ & Period & Minimum inter-arrival time (ticks) \\
$D_i \in \N^+$ & Deadline & Relative deadline ($D_i = T_i$ for implicit deadlines) \\
$C_i^{\mathrm{lo}} \in \N^+$ & Normal budget & Execution budget in normal mode \\
$C_i^{\mathrm{hi}} \in \N^+$ & HI execution budget & Execution budget once fully degraded ($C_i^{\mathrm{hi}} \geq C_i^{\mathrm{lo}}$) \\
$\mathrm{BCET}_i \in \N^+$ & Best-case execution time & $\mathrm{BCET}_i \leq C_i^{\mathrm{lo}}$ \\
\bottomrule
\end{tabular}
\end{table}

\subsection{Severity and Per-Level Budgets}
\label{sec:severity-budgets}

Every degradation level has its own execution budget. Rather than name each
level's budget independently, we index it by a \emph{severity}
$\chi \in [0, 1]$ and interpolate:

\[
C_i(\chi) = \max\bigl(C_i^{\mathrm{lo}},\; C_i^{\mathrm{lo}} + \chi \cdot (C_i^{\mathrm{hi}} - C_i^{\mathrm{lo}})\bigr)
\]

so that $C_i(0) = C_i^{\mathrm{lo}}$, $C_i(1) = C_i^{\mathrm{hi}}$, and $\chi$
between the two endpoints gives an intermediate budget. The outer $\max$ is not
decorative: it is what guarantees $C_i(\chi) \geq C_i^{\mathrm{lo}}$ for every
$\chi$, which in turn is what makes the response times derived from it never
earlier than the normal-mode response time (\Cref{prop:c} below) -- rounding
must never be allowed to violate it. Linear interpolation is a design choice,
not a derivation: every property in this document holds for any interpolation
that is monotone in $\chi$ and never below $C_i^{\mathrm{lo}}$, so a
better-motivated member of that family is an open question, not a constraint
this analysis depends on.

\subsection{Criticality Assignment}

Each task is assigned a criticality level:
\[
\kappa_i \in \{\mathrm{HI}, \mathrm{LO}\}
\]

Let $\tau^{\mathrm{HI}} = \{\tau_i \in \tau \mid \kappa_i = \mathrm{HI}\}$ and
$\tau^{\mathrm{LO}} = \{\tau_i \in \tau \mid \kappa_i = \mathrm{LO}\}$ denote the
HI-criticality and LO-criticality subsets, respectively.

\subsection{Utilisation}

The utilisation of task $\tau_i$ under different criticality modes is:

\[
U_i^{\mathrm{lo}} = \frac{C_i^{\mathrm{lo}}}{T_i}, \quad
U_i^{\mathrm{hi}} = \begin{cases}
    \frac{C_i^{\mathrm{hi}}}{T_i} & \text{if } \kappa_i = \mathrm{HI} \\
    0 & \text{if } \kappa_i = \mathrm{LO}
\end{cases}
\]

The total LO-criticality utilisation is $U = \sum_{i=1}^n U_i^{\mathrm{lo}}$.
The HI-criticality utilisation factor (CF) relates aggregate utilisations:
\[
U^{\mathrm{hi}} = \sum_{\tau_i \in \tau^{\mathrm{HI}}} U_i^{\mathrm{hi}}, \quad
\text{where } U^{\mathrm{hi}} = CP \cdot CF \cdot U
\]
Here $CP = |\tau^{\mathrm{HI}}| / n$ is the criticality proportion.

\section{Response Time Analysis}
\label{sec:rta}

This section presents the standard fixed-priority response-time analysis, then
generalises it to a severity $\chi$. It is placed before the degradation model
(\Cref{sec:multilevel}) because that model's trigger points and safety argument
are both stated in terms of the response times defined here.

\subsection{Normal-Mode Response Time}
\label{sec:r-lo}

For fixed-priority preemptive scheduling, the response time of task $\tau_i$
under normal-mode budgets is the least fixed point of
\[
R_i^{\mathrm{lo}} = C_i^{\mathrm{lo}} + \sum_{\tau_j \in \mathrm{hp}(i)} \left\lceil \frac{R_i^{\mathrm{lo}}}{T_j} \right\rceil C_j^{\mathrm{lo}}
\]
where $\mathrm{hp}(i)$ is the set of tasks with higher priority than $\tau_i$.
\textbf{Requirement R1}, that a task set meets every deadline when every job
complies with its LO budget, holds exactly when $R_i^{\mathrm{lo}} \leq D_i$ for
every $\tau_i \in \tau$.

\subsection{Classic Two-Level Schedulability (AMC-rtb)}
\label{sec:amc-rtb}

The design-time question "can this task set tolerate full degradation at
all?" is answered independently of any run-time trigger, by the AMC-rtb test.
For a HI-criticality task $\tau_i$, the response time once every LO-criticality
task has been abandoned is the least fixed point of
\[
R_i^{\mathrm{hi}} = C_i^{\mathrm{hi}}
  + \sum_{\tau_j \in \mathrm{hp}_{\mathrm{HI}}(i)} \left\lceil \frac{R_i^{\mathrm{hi}}}{T_j} \right\rceil C_j^{\mathrm{hi}}
  + \sum_{\tau_k \in \mathrm{hp}_{\mathrm{LO}}(i)} \min\!\left( \left\lceil \frac{R_i^{\mathrm{hi}}}{T_k} \right\rceil,\; \left\lceil \frac{R_i^{\mathrm{lo}}}{T_k} \right\rceil + 1 \right) C_k^{\mathrm{lo}}
\]
where $\mathrm{hp}_{\mathrm{HI}}(i)$ and $\mathrm{hp}_{\mathrm{LO}}(i)$ are the
higher-priority HI-criticality and LO-criticality tasks, respectively. The
$\min(\cdot,\cdot)$ term bounds the interference from a higher-priority
LO-criticality task by the number of its jobs that can still be in flight when
degradation begins -- at most one more than it would have released by
$R_i^{\mathrm{lo}}$ -- since every LO-criticality task is abandoned on release
once fully degraded. A task set is AMC-rtb schedulable if $R_i^{\mathrm{hi}}
\leq D_i$ for every $\tau_i \in \tau^{\mathrm{HI}}$ and $R_i^{\mathrm{lo}} \leq
D_i$ for every $\tau_i \in \tau$.

This is the two-level baseline this work compares against, and the degenerate
case ($k = 2$, every LO-criticality task abandoned at the single trigger) that
the general model in the next section must reduce to exactly.

\subsection{Severity-Indexed Response Time}
\label{sec:severity-response}

Generalise \Cref{sec:r-lo} by charging \emph{every} task at severity $\chi$
rather than only at $C_i^{\mathrm{lo}}$:
\[
R_i(\chi) = C_i(\chi) + \sum_{\tau_j \in \mathrm{hp}(i)} \left\lceil \frac{R_i(\chi)}{T_j} \right\rceil C_j(\chi)
\]
using the per-task budgets of \Cref{sec:severity-budgets}. If this recurrence
exceeds $D_i$ before converging, level $\chi$ is unreachable for task $i$ and
$R_i(\chi) = \infty$ -- a threshold that never fires, rather than the
partially-iterated value the recurrence had reached, which would depend on
where iteration happened to stop and break the monotonicity property below.

$R_i(\chi)$ answers a different question from $R_i^{\mathrm{hi}}$ in
\Cref{sec:amc-rtb}: it is the response time \emph{if every task, including
every LO-criticality one, were uniformly behaving at severity $\chi$ with none
yet abandoned}. That is deliberately more pessimistic than the true response
time once a drop policy is in effect, and pessimism in this direction is what
makes $R_i(\chi)$ safe to use as a run-time trigger: reaching it is evidence
that observed behaviour is \emph{at least} that severe, never a false alarm.
$R_i^{\mathrm{hi}}$, by contrast, models the true operating response time under
the specific policy "abandon every LO-criticality task", and is a design-time
schedulability test, not something monitored at run time.

Three properties hold for every task set, and are what the multi-level
scheme's safety argument (\Cref{sec:admissibility}) rests on. They are checked
directly, not merely argued, in \texttt{tests/scheduling/test\_severity.py}.

\begin{description}
\item[(A)] $R_i(0) = R_i^{\mathrm{lo}}$ exactly -- a level at severity $0$
  triggers at exactly the classic single AMC trigger point.
\item[(B)] $\chi_a \leq \chi_b \implies R_i(\chi_a) \leq R_i(\chi_b)$ -- the
  severity-to-response-time map is monotone, so a ladder of increasing
  severities produces a ladder of non-decreasing triggers, given the
  $\infty$-saturation above.
\item[(C)] $R_i(\chi) \geq R_i^{\mathrm{lo}}$ for every $\chi$ -- no severity
  level ever triggers earlier than the classic AMC trigger.
\end{description}
\label{prop:c}

Property (C) is what an earlier fraction-based design ($\chi \cdot
R_i^{\mathrm{lo}}$ for $\chi < 1$, triggering \emph{below} the classic
threshold rather than at or above it) violated by construction: measured with
no HI-criticality behaviour present anywhere, such a threshold at $0.9 \cdot
R_i^{\mathrm{lo}}$ fired 14{,}118 times across twelve task sets, abandoning 356
jobs with no fault present. Charging every task's own budget, rather than
scaling the deadline, is what avoids this.

\section{Multi-Level Degradation Model}
\label{sec:multilevel}

\subsection{Degradation Levels}

The system has $k \geq 2$ levels $L_0, L_1, \dots, L_{k-1}$, each assigned a
severity $0 = \chi_1 \leq \chi_2 \leq \dots \leq \chi_{k-1} \leq 1$ (level
$L_0$ has no severity; it is the absence of any triggered level). Pinning
$\chi_1 = 0$ is what makes $L_1$ trigger at exactly the classic AMC point
(property A) and is required by the safety argument in
\Cref{sec:admissibility}; $\chi_{k-1}$ need not be $1$, though setting it so
recovers the classic two-level scheme's full protection at the deepest level.

\begin{table}[h]
\centering
\caption{Degradation Level Structure}
\begin{tabular}{cccc}
\toprule
Level & Severity & Trigger Condition & LO Task Handling \\
\midrule
$L_0$ & -- & None & All tasks run at $C_i^{\mathrm{lo}}$ \\
$L_1$ & $\chi_1 = 0$ & HI task reaches $R_i(0) = R_i^{\mathrm{lo}}$ & Drop set $S_1$ (Section~\ref{sec:admissibility}) \\
$L_2$ & $\chi_2$ & HI task reaches $R_i(\chi_2)$ & Drop set $S_2 \supseteq S_1$ \\
$\vdots$ & $\vdots$ & $\vdots$ & $\vdots$ \\
$L_{k-1}$ & $\chi_{k-1}$ & HI task reaches $R_i(\chi_{k-1})$ & Drop set $S_{k-1} \supseteq \dots \supseteq S_1$ \\
\bottomrule
\end{tabular}
\end{table}

Each level's operating budget for every task is $C_i(\chi_x)$
(\Cref{sec:severity-budgets}): a task entering $L_x$ because severity $\chi_x$
was reached is granted exactly the execution allowance that severity
justifies, not immediately the full $C_i^{\mathrm{hi}}$ -- that is what makes
the degradation graded rather than binary. One consequence is worth stating
plainly: with $\chi_1 = 0$, $L_1$'s operating budget equals $L_0$'s. Entering
$L_1$ changes nothing about what any task is allowed to execute; it only marks
that an overrun has been confirmed somewhere; and it must still ask whether any
LO-criticality task needs to be shed to keep HI tasks safe at that (unchanged)
budget, which \Cref{sec:admissibility} answers per level. This is the same
principle Theorem 1 of the safety proof states for $L_0$ itself: protection
begins no later than the classic trigger, never earlier.

\subsection{Mode Transition Protocol}

The system transitions between degradation levels according to:
\begin{align*}
L_x \to L_{x+1} &\quad \text{when } \exists\, \tau_i \in \tau^{\mathrm{HI}}: \text{active job reaches } R_i(\chi_{x+1}) \\
L_x \to L_y &\quad y < x \text{ on exit (direct or cascade, per protocol design)}
\end{align*}
The trigger point for the deepest level, $R_i(\chi_{k-1})$, is written
$R_{\mathrm{trigger}}$ where the level count is not otherwise relevant; with
$\chi_{k-1} = 1$ it recovers the classic AMC trigger for full degradation.

\section{Admissible Drop Sets and Safety}
\label{sec:admissibility}

Retaining LO-criticality work at an intermediate level necessarily means
absorbing interference that abandoning it would have avoided -- so the drop
set $S_x$ at each level is not a free design choice, it is an obligation
discharged by response-time analysis at design time.

\textbf{Definition (admissible drop set).} $S_x \subseteq \tau^{\mathrm{LO}}$
is \emph{admissible} at severity $\chi_x$ if, charging every task
$C_i(\chi_x)$, treating tasks in $S_x$ as absent, and every retained
LO-criticality task at $C_i^{\mathrm{lo}}$:
\begin{enumerate}
\item every HI-criticality task $\tau_i$ has response time $R_i \leq D_i$, and
\item every retained LO-criticality task $\tau_j \notin S_x$ has response time
  $R_j \leq D_j$.
\end{enumerate}

Clause 2 is what makes LDM (LO jobs that miss their deadline in degraded mode)
zero by construction for any retained task: there is no
JNE-against-LDM trade-off to measure, only JNE against how much of the
severity range is covered without shedding. An admissible set always exists
whenever the task set passes AMC-rtb (\Cref{sec:amc-rtb}): $S_x = \tau^{\mathrm{LO}}$
reduces clause 1 to the AMC-rtb HI-mode test and vacates clause 2, recovering
the classic two-level scheme exactly. So the question a drop policy answers is
never \emph{whether} a level is feasible, only how little must be shed --
which is what makes the drop set an optimisation rather than a guess, and is
addressed separately (drop-strategy research, out of scope for this document).

Measured over 50 AMC-rtb-schedulable task sets ($n=12$, $U=0.7$, $CF=2$),
minimal admissible sets shed markedly less than the classic scheme's blanket
100\%:

\begin{table}[h]
\centering
\caption{Minimal admissible LO-criticality tasks shed, by severity}
\begin{tabular}{lc}
\toprule
Severity $\chi$ & LO tasks shed \\
\midrule
$\leq 0.25$ & 0.0\% \\
0.50 & 10.0\% \\
0.75 & 33.7\% \\
1.00 & 56.3\% \\
\midrule
classic two-level (single trigger) & 100\% \\
\bottomrule
\end{tabular}
\end{table}

\textbf{Nesting.} $S_1 \subseteq S_2 \subseteq \dots \subseteq S_{k-1}$ is
required so that a level transition only ever adds to the drop set, keeping
the run-time decision incremental. This holds automatically for any drop
strategy that is a pure function of the task set and the remaining candidates
(such greedy strategies shed the same sequence of tasks regardless of
severity, so each level's set is a prefix of one sequence, and prefixes nest)
-- it is not an extra assumption on top of the definition above.

\textbf{Scope of this analysis.} The response-time bound above treats a
shed task's interference as ceasing the instant its level is entered. The
simulator, like the classic scheme, instead lets a job already released
continue to completion and abandons only that task's \emph{future} releases --
so a task shed at the transition instant can still contribute the interference
of one job already in flight, exactly the quantity the $\min(\cdot,\cdot)$ term
in \Cref{sec:amc-rtb} bounds for the classic two-level case. The admissibility
check in this section does not yet include that term for a general drop set;
doing so is a direct generalisation of \Cref{sec:amc-rtb}'s bound and is
necessary before this analysis is relied on as a safety guarantee, rather than
as the sound-in-the-limit design it currently is.

\section{Job Execution Model}

\subsection{Execution Time Distribution}

Each job's execution time is drawn at release:
\[
E_i \sim
\begin{cases}
\mathrm{Uniform}(\mathrm{BCET}_i, C_i^{\mathrm{lo}}) & \text{if } \kappa_i = \mathrm{LO} \\
\mathrm{Uniform}(\mathrm{BCET}_i, C_i^{\mathrm{lo}}) & \text{if } \kappa_i = \mathrm{HI}, \text{normal behaviour} \\
\mathrm{Uniform}(C_i^{\mathrm{lo}}, C_i^{\mathrm{hi}}) & \text{if } \kappa_i = \mathrm{HI}, \text{HI behaviour}
\end{cases}
\]
HI-criticality behaviour occurs with probability $FP = 1/N$ per job.

\subsection{Execution Budget Enforcement}

During execution, a job is constrained by the budget of the level the system
is in \emph{when that job is released}:
\[
\text{Actual execution} = \min(E_i, C_i(\chi_x))
\]
where $\chi_x$ is the severity of the current level $L_x$ ($\chi_1 = 0$ at the
shallowest degraded level, up to $\chi_{k-1}$ at the deepest). This is a
budget per level, not per mode: at $k = 2$ it recovers $B_i = C_i^{\mathrm{lo}}$
in normal mode and $B_i = C_i^{\mathrm{hi}}$ once degraded.

\section{Notation Summary}

\begin{table}[h]
\centering
\caption{Notation Reference}
\begin{tabular}{ll}
\toprule
Symbol & Meaning \\
\midrule
$\tau$ & Task set \\
$n$ & Number of tasks \\
$k$ & Number of degradation levels ($k \geq 2$) \\
$\kappa_i$ & Criticality of task $i$ ($\mathrm{HI}$ or $\mathrm{LO}$) \\
$T_i, D_i$ & Period and deadline of task $i$ \\
$C_i^{\mathrm{lo}}, C_i^{\mathrm{hi}}$ & Normal and HI-criticality execution budgets \\
$\chi \in [0,1]$ & Severity; $0$ = normal, $1$ = fully degraded \\
$C_i(\chi)$ & Execution budget for task $i$ at severity $\chi$ \\
$\mathrm{BCET}_i$ & Best-case execution time \\
$U_i^{\mathrm{lo}}, U_i^{\mathrm{hi}}$ & Utilisation in LO/HI modes \\
$R_i^{\mathrm{lo}}$ & Normal-mode response time (Section~\ref{sec:r-lo}) \\
$R_i^{\mathrm{hi}}$ & Classic two-level HI-mode response time (Section~\ref{sec:amc-rtb}) \\
$R_i(\chi)$ & Severity-indexed response time / run-time trigger threshold (Section~\ref{sec:severity-response}) \\
$R_{\mathrm{trigger}}$ & $R_i(\chi_{k-1})$, the deepest level's trigger \\
$\chi_x$ & Severity assigned to level $L_x$ \\
$S_x$ & Drop set at level $L_x$ (Section~\ref{sec:admissibility}) \\
$FP = 1/N$ & HI-behaviour probability per job \\
$\tau^{\mathrm{HI}}, \tau^{\mathrm{LO}}$ & HI/LO criticality subsets \\
$L_x$ & Degradation level $x$ ($0 \leq x < k$) \\
\bottomrule
\end{tabular}
\end{table}

\bigskip
\noindent\textit{References are drawn first from the project-specific bibliography
(\texttt{project\_bib.bib}), then from the master Overleaf bibliography
(\texttt{references.bib}).}

\bibliographystyle{plain}
\bibliography{project_bib,references}

\end{document}


\section{Evaluation Methodology}

See the evaluation methodology document for details on:
\begin{itemize}
    \item Task set generation using DRS algorithm
    \item Metrics (NiD, TiD, JNE, LDM, HDM)
    \item Objective function formulation
\end{itemize}

\documentclass{article}
\usepackage{amsmath, amssymb, amsthm}
\usepackage{booktabs}
\usepackage{geometry}
\usepackage{hyperref}

\geometry{margin=1in}
\hypersetup{colorlinks=true, urlcolor=blue}

% Shared macros for the multi-mode mixed-criticality scheduling papers.
% % Shared macros for the multi-mode mixed-criticality scheduling papers.
% \input{macros} from any document in this directory rather than redefining
% these locally -- a macro defined in one paper's preamble and silently
% missing from a sibling's is exactly the class of error that put \round and
% \E undefined in evaluation_methodology.tex.

\newcommand{\N}{\mathbb{N}}
\newcommand{\R}{\mathbb{R}}
\newcommand{\Z}{\mathbb{Z}}
\newcommand{\E}{\mathbb{E}}
\newcommand{\Var}{\mathrm{Var}}
\newcommand{\Prob}{\mathbb{P}}
\newcommand{\Ind}{\mathbf{1}}
\newcommand{\ceil}[1]{\lceil #1 \rceil}
\newcommand{\floor}[1]{\lfloor #1 \rfloor}
\newcommand{\round}[1]{\lfloor #1 \rceil}
 from any document in this directory rather than redefining
% these locally -- a macro defined in one paper's preamble and silently
% missing from a sibling's is exactly the class of error that put \round and
% \E undefined in evaluation_methodology.tex.

\newcommand{\N}{\mathbb{N}}
\newcommand{\R}{\mathbb{R}}
\newcommand{\Z}{\mathbb{Z}}
\newcommand{\E}{\mathbb{E}}
\newcommand{\Var}{\mathrm{Var}}
\newcommand{\Prob}{\mathbb{P}}
\newcommand{\Ind}{\mathbf{1}}
\newcommand{\ceil}[1]{\lceil #1 \rceil}
\newcommand{\floor}[1]{\lfloor #1 \rfloor}
\newcommand{\round}[1]{\lfloor #1 \rceil}


\title{Evaluation Methodology for Multi-Mode Mixed-Criticality Scheduling}
\author{}
\date{}

\begin{document}

\maketitle

\section{Task Set Generation}

Following DRS (RTSS 2020)~\cite{1a02ee59d96a4af497808eeb685b574c}, utilisation vectors are generated via the Dirichlet-Rescale algorithm.

\subsection{DRS Algorithm}

Given $n$ tasks and target utilisation $U$, the DRS algorithm produces $\mathbf{u} = (u_1, \dots, u_n)$ such that:
\[
\sum_{i=1}^n u_i = U, \quad 0 \leq u_i \leq 1
\]

The algorithm operates via:
\begin{enumerate}
    \item Draw $\mathbf{x} \sim \mathrm{Dirichlet}(1, \dots, 1)$ (uniform on simplex)
    \item Scale: $u_i = U \cdot x_i$
    \item Apply constraints via rescaling to standard simplex
\end{enumerate}

\subsection{Task Parameter Generation}

For each task $\tau_i$, parameters are generated as follows:

\begin{align*}
T_i &\sim \mathrm{LogUniform}(\log(T_{\min}), \log(T_{\max})) \\
D_i &= T_i \\
C_i^{(0)} &= \round(u_i \cdot T_i) \\
\mathrm{BCET}_i &\sim \mathrm{Uniform}(BCET_{\mathrm{min}} \cdot C_i^{(0)}, C_i^{(0)}) \\
C_i^{(k-1)} &= \min(\round(CF \cdot C_i^{(0)}), T_i)
\end{align*}

where:
\begin{itemize}
    \item $T_{\min}$ and $T_{\max}$ are the minimum and maximum period bounds
    \item $BCET_{\mathrm{min}}$ is the minimum BCET fraction (typically $0.8$)
    \item $CF$ is the criticality factor relating normal to degraded budgets
\end{itemize}

\section{Metrics}

\subsection{Degradation Metrics}

\begin{align*}
\textbf{NiD} &\quad \text{Number of times degraded mode is entered} \\
\textbf{TiD} &\quad \text{Fraction of time spent in degraded mode} \\
\textbf{JNE} &\quad \text{Jobs Not Executed (LO jobs dropped)} \\
\textbf{LDM} &\quad \text{Late Degraded Mode (LO jobs that miss deadline in degraded mode)} \\
\textbf{HDM} &\quad \text{High Degradation Mode (HI jobs missing deadline)}
\end{align*}

\subsection{Objective Function}

For $k$-level scheduling, we define the objective function:

\[
\Phi = \alpha \cdot \E[\mathrm{JNE}] + \beta \cdot \E[\mathrm{TiD}] + \gamma \cdot \E[\mathrm{WastedCPU}] + \delta \cdot \E[\mathrm{NiD}]
\]

where:
\begin{itemize}
    \item $\mathrm{JNE}$: Jobs Not Executed (LO jobs dropped)
    \item $\mathrm{TiD}$: Fraction of time spent in degraded mode
    \item $\mathrm{WastedCPU}$: CPU cycles spent on abandoned LO work
    \item $\mathrm{NiD}$: Number of times degraded mode is entered (penalizes frequent mode changes)
\end{itemize}
The coefficients $\alpha, \beta, \gamma, \delta$ are weighting parameters that can be tuned based on
application requirements.

\section{Simulation Parameters}

Common parameter values used in evaluation:

\begin{table}[h]
\centering
\caption{Default Simulation Parameters}
\begin{tabular}{lll}
\toprule
Parameter & Default Value & Description \\
\midrule
$T_{\min}$ & 100 & Minimum task period (ticks) \\
$T_{\max}$ & 10000 & Maximum task period (ticks) \\
$BCET_{\mathrm{min}}$ & 0.8 & Minimum BCET as fraction of $C_i^{(0)}$ \\
$CF$ & See Section~\ref{sec:parameters} & Criticality factor \\
$k$ & 2 & Number of degradation levels \\
\bottomrule
\end{tabular}
\end{table}

\section{Notation Summary}

\begin{table}[h]
\centering
\caption{Notation Reference}
\begin{tabular}{ll}
\toprule
Symbol & Meaning \\
\midrule
$n$ & Number of tasks \\
$U$ & Target utilisation \\
$\mathbf{u}$ & Utilisation vector \\
$T_i$ & Period of task $i$ \\
$D_i$ & Deadline of task $i$ \\
$C_i^{(0)}$ & Normal execution budget \\
$C_i^{(k-1)}$ & HI execution budget \\
$\mathrm{BCET}_i$ & Best-case execution time \\
$U_i^{\mathrm{lo}}, U_i^{\mathrm{hi}}$ & Utilisation in LO/HI modes \\
$FP = 1/N$ & HI-behaviour probability per job \\
$CF$ & Criticality factor \\
$k$ & Number of degradation levels \\
$\tau^{\mathrm{HI}}, \tau^{\mathrm{LO}}$ & HI/LO criticality subsets \\
\bottomrule
\end{tabular}
\end{table}

\section{References}

\bibliographystyle{plain}
\bibliography{project_bib,references}

\end{document}


\end{document}
